\documentclass[10pt,a4paper]{amsart}
\usepackage[margin=19mm]{geometry}
\usepackage{amsmath,amssymb,mathtools}
\usepackage[scr=rsfs,cal=euler]{mathalfa}
\usepackage{xfrac}
\usepackage[unicode,hidelinks,bookmarks=false]{hyperref}
\newcommand{\ie}{i.e.\ }
\newcommand{\cf}{cf.\ }
\newcommand{\resp}{respectively\ }
\newcommand{\Subsection}{\subsection}
\providecommand{\nobreakdash}{\nobreak\hbox{-}}
\setlength{\emergencystretch}{2em}
\multlinegap0pt
\usepackage{bm}
\usepackage{bbm}
\usepackage{dsfont}
\usepackage{xspace}
\usepackage{cancel}
\usepackage{mathtools}
\usepackage{amsthm}
\makeatletter
\@ifundefined{equation}{\newtheorem{equation}{Equation}}{}
\@ifundefined{Proposition}{\newtheorem{Proposition}[equation]{Proposition}}{}
\@ifundefined{Theorem}{\newtheorem{Theorem}[equation]{Theorem}}{}
\makeatother
\def\Z{\mathbb{Z}}
\renewcommand{\vec}[1]{{\bf #1}}
\title[The conjugate orbit of a unitary operator]{The conjugate orbit of a unitary operator}
\author{Javad Mashreghi, Marek Ptak,, William T. Ross}
\date{Source: arXiv:2602.19337v1 (CC BY 4.0)}
\begin{document}
\maketitle

\noindent\emph{Source and licence.} The conjugate orbit of a unitary operator. Version arXiv:2602.19337v1 (CC BY 4.0), distributed under the Creative Commons Attribution 4.0 International licence, as recorded in the arXiv licence metadata for this exact version and shown on the abstract page https://arxiv.org/abs/2602.19337v1. Full rights notice: licenses/arxiv-2602.19337-CC-BY-4.0.txt.\par\smallskip

\noindent\textit{Editorial scope.} Counted unit: the Theorem whose source label is \texttt{0934t4rhhHHCccvvJ} in the article (source0.tex lines 882--889, source label \texttt{0934t4rhhHHCccvvJ}), delivered together with the article's own complete original proof of it (source0.tex lines 891--933, one \texttt{proof} environment of 43 lines). No mathematics is written, repaired or re-derived: statement and proof are the article's own text line for line. Required context retained uncounted: OShfgfdsfa (lines 868--874). Every definition, display and label of those lines is retained unchanged. Omitted from this package and not redistributed: 1--867, 875--881, 890--890, 934--2356. Nothing inside a retained excerpt is omitted or altered: every retained body is delivered exactly as the article's lines are written, so no omission receipt of any kind accompanies this package. Citations: the retained excerpts carry no citation command at all, so no bibliography entry of the article is transcribed and no reference is redistributed. Reference adaptation: none was needed and none was made; every reference inside the delivered unit resolves inside it, so no unresolved reference or citation is produced. Printed slips kept literal: no printed word, symbol, hypothesis or number of the counted statement or of its proof was altered. Layout and class: the article's own class is \texttt{dis} with packages including amsmath, amssymb, amsthm, amsfonts, enumerate, color, comment, mathrsfs, parskip, xcolor, mathtools, palatino, graphicx, enumitem; the wrapper is the lane's pinned \texttt{amsart} editorial wrapper at 10pt on A4, which restates the theorem-like environments the retained text needs and adds the article's own macro definitions that the retained text needs (\texttt{\textbackslash Z}, \texttt{\textbackslash vec}), with extra wrapper packages (bm, bbm, dsfont, xspace, cancel, mathtools). The delivered numbering is the unit's own. Assets: no retained span contains a figure environment, a graphics-inclusion command, a PDF-page inclusion, a table or a TikZ picture; no image file is copied into this package, the package declares no asset dependency and no image file is redistributed with it. Licence: version arXiv:2602.19337v1 of this article is distributed under the Creative Commons Attribution 4.0 International licence, as recorded in the arXiv licence metadata for this exact version and shown on the abstract page https://arxiv.org/abs/2602.19337v1. A full rights notice is delivered at licenses/arxiv-2602.19337-CC-BY-4.0.txt. Characters: every retained excerpt is pure ASCII, so the delivered engine, XeLaTeX with its default Latin Modern text font, reports no missing character.
\begin{Proposition}\label{OShfgfdsfa}
For a map $C$ on $\ell^2(\Z)$, the following are equivalent. 
\begin{enumerate}
\item $C$ is a conjugation; 
\item There is a doubly infinite unitary matrix with $V = [V_{i j}]_{i, j \in \Z}$ with $V = V^{t}$ and $C = V J$. In other words, $C \vec{x} = V \overline{\vec{x}}$.
\end{enumerate}
\end{Proposition}

\begin{Theorem}\label{0934t4rhhHHCccvvJ}
For a unitary operator $W$ on $\ell^2(\Z)$, the following are equivalent. 
\begin{enumerate}
\item $W \in \mathfrak{O}_c(M)$; 
\item $W$ is a unitary shift with $W \vec{v}_{j} = \vec{v}_{j + 1}$, where $V = [\cdots| \vec{v}_{-1}|\vec{v}_0|\vec{v}_1|\cdots]$ is a doubly infinite unitary matrix with $V^{t} = V$. Moreover, $W = (VJ) M (VJ)$.
\end{enumerate}
%More specifically,  if  $V$ is a doubly infinite unitary matrix with $V^{t} = V$ and  $C= V J$ is a conjugation on $\ell^2(\Z)$, then $W = C M C$ is a unitary shift with respect to the orthonormal basis $(V \vec{e}_n)_{n \in \Z}$ for $\ell^2(\Z)$. Moreover, if $V$ is any doubly infinite unitary matrix with $V = V^{t}$, then there is a conjugation $C$ on $\ell^2(\Z)$ such that $C M C (V \vec{e}_n)= V \vec{e}_{n + 1}$, that is, $W = C M C$ is a unitary shift on $\ell^2(\Z)$.
\end{Theorem}

\begin{proof}
$(a) \Longrightarrow (b)$: If $W \in \mathfrak{O}_c(M)$ then, by definition of the conjugate orbit, there is a conjugation $C$ on $\ell^2(\Z)$ with $C M C = W$. By Proposition \ref{OShfgfdsfa}, $C = V J$, where  $V$ is a doubly infinite self transpose unitary matrix. Then, with respect to the orthonormal basis  $(V \vec{e}_{n})_{n \in \Z}$ for $\ell^2(\Z)$, which are just the columns of the matrix  $V$, we have
\begin{align*}
W (V \vec{e}_n) & = 
C M C (V \vec{e}_n)\\
 & = C M C V J \vec{e}_n\\
& = C M C^2 \vec{e}_n \\
&= C M \vec{e}_n \\
&= C \vec{e}_{n + 1} \\
& = V J \vec{e}_{n + 1} \\
& = V \vec{e}_{n + 1}.
\end{align*}
Thus, $W$ is a unitary shift of the desired form.

$(b) \Longrightarrow (a)$: 
Suppose that $W$ is a unitary shift with corresponding column vectors arising from a doubly infinite unitary matrix  $V$ with $V = V^{t}$ (equivalently, $J V J = V^{*}$),  in that $W \vec{v}_n = \vec{v}_{n + 1}$ with $\vec{v}_{n} = V \vec{e}_n$ for all $n \in \Z$.  Then $C = V J$ is a conjugation on $\ell^2(\Z)$. 
%Indeed, $C$ is antilinear and isometric with 
%$$C^2 = V J V J = V \overline{V} = V V^{*} = I.$$ 
Moreover,
\begin{align*}
(CMC) \vec{v}_n & = (C M C) V \vec{e}_n\\
& = C M C C J \vec{e}_n\\
&= C M C C \vec{e}_{n}\\
& = C M \vec{e}_{n}\\
& = C \vec{e}_{n + 1}\\
& = C J \vec{e}_{n + 1}\\
& = V \vec{e}_{n + 1}\\
& = \vec{v}_{n + 1}\\
& = W \vec{v}_n.
\end{align*}
The above says that the two unitary operators $W$ and $CMC$ on $\ell^2(\Z)$ agree on the orthonormal basis $(\vec{v}_n)_{n = -\infty}^{\infty}$ (and hence everywhere) and so $W \in \mathfrak{O}_c(M)$. 

To verify that $W = (VJ) M (VJ)$, observe that since $V^{t} = V$, we see that $V \overline{\vec{v}_j} = \vec{e}_j$. Thus, 
\begin{align*}
(VJ) M (VJ) \vec{v}_j & = VJ M V \overline{\vec{v}_j}\\
& = VJ M \vec{e}_j\\
& = VJ \vec{e}_{j + 1}\\
& = V \vec{e}_{j + 1}\\
& = \vec{v}_{j + 1}\\
& = W \vec{v}_{j},
\end{align*}
which proves that $W = (V J) M (VJ)$.
\end{proof}

\end{document}
